% Môn: Vật lý
% Lớp: 12
% Chương: Chương I. Vật lí nhiệt
% Bài: L12C1 Bài 4. Nhiệt dung riêng · Bài toán công suất, hiệu suất và thời gian hóa hơi
% ===== Câu 1 =====
% ID: L12C1B4-88-DS
% Mức: NB
\dangbt{Bài toán công suất, hiệu suất và thời gian hóa hơi}
\begin{ex}
% ID: L12C1B4-88-DS
% Mức: NB
Cùng chất, cùng khối lượng và hiệu suất.
\choiceTF
{\True t tỉ lệ nghịch P.}
{\True Tăng $P$ gấp đôi làm t giảm một nửa.}
{Nếu P=0 thì vẫn hóa hơi trong thời gian hữu hạn do công thức.}
{\True Công suất hao phí là (1-H)P.}
\loigiai{
\begin{itemchoice}
\item (Đúng) t tỉ lệ nghịch P.
\item (Đúng) Tăng $P$ gấp đôi làm t giảm một nửa.
\item (Sai) Nếu P=0 thì vẫn hóa hơi trong thời gian hữu hạn do công thức.
\item (Đúng) Công suất hao phí là (1-H)P.
\end{itemchoice}
}
\end{ex}

% ===== Câu 2 =====
% ID: L12C1B4-89-TL
% Mức: NB
\begin{bt}
% ID: L12C1B4-89-TL
% Mức: NB
Chứng minh thời gian hóa hơi t tỉ lệ thuận m và nghịch với $P$ khi $L$, $H$ không đổi. Nêu cách rút ngắn thời gian.
\loigiai{
Từ HPt=Lm suy ra t=Lm/(HP). Với $L$, $H$ cố định, t∝m và t∝1/P. Có thể giảm m, tăng công suất hữu ích HP hoặc giảm hao phí để rút ngắn thời gian.
}
\end{bt}

% ===== Câu 3 =====
% ID: L12C1B4-90-TLN
% Mức: NB
\begin{ex}
% ID: L12C1B4-90-TLN
% Mức: NB
Nguồn 1000 $W$ hoạt động 10 phút, hiệu suất 80%, L=2,0.10^ $6\,\text{J}$ /kg. Khối lượng hóa hơi theo gam là bao nhiêu? Chỉ ghi số.
\shortans{240}
\loigiai{
$m=HPt/L=0,8.1000.600/(2,0.10^6)=0,24 kg=240 g.$
}
\end{ex}

% ===== Câu 4 =====
% ID: L12C1B4-91-TN
% Mức: NB
\begin{ex}
% ID: L12C1B4-91-TN
% Mức: NB
Nguồn công suất 1000 $W$, hiệu suất 80%, hóa hơi $0,12\,\text{kg}$ chất có L=2.10^ $6\,\text{J}$ /kg ở nhiệt độ sôi. Thời gian là
\choice
{$30\,\text{s}$}
{\True $300\,\text{s}$}
{$3000\,\text{s}$}
{300 phút}
\loigiai{
H.P.t=Lm nên t=Lm/(HP)= $300\,\text{s}$.
}
\end{ex}

% ===== Câu 5 =====
% ID: L12C1B4-92-DS
% Mức: TH
\begin{ex}
% ID: L12C1B4-92-DS
% Mức: TH
Một thiết bị có P=1500 $W$, H=80%.
\choiceTF
{\True Công suất hữu ích là $1200\,\text{W}$.}
{\True Trong 10 phút, năng lượng hữu ích là 720 kJ.}
{\True Năng lượng tiêu thụ trong 10 phút là 900 kJ.}
{Năng lượng hữu ích lớn hơn năng lượng tiêu thụ.}
\loigiai{
\begin{itemchoice}
\item (Đúng) Công suất hữu ích là $1200\,\text{W}$.
\item (Đúng) Trong 10 phút, năng lượng hữu ích là 720 kJ.
\item (Đúng) Năng lượng tiêu thụ trong 10 phút là 900 kJ.
\item (Sai) Năng lượng hữu ích lớn hơn năng lượng tiêu thụ.
\end{itemchoice}
}
\end{ex}

% ===== Câu 6 =====
% ID: L12C1B4-93-TL
% Mức: TH
\begin{bt}
% ID: L12C1B4-93-TL
% Mức: TH
Một bếp 1500 $W$, hiệu suất 80% hóa hơi $0,45\,\text{kg}$ nước ở 100 °C. Lấy L=2,3.10^ $6\,\text{J}$ /kg. Tính thời gian.
\loigiai{
t=Lm/(HP)=2,3.10^6.0,45/(0,8.1500)= $862,5\,\text{s}$, khoảng 14,4 phút.
}
\end{bt}

% ===== Câu 7 =====
% ID: L12C1B4-94-TLN
% Mức: NB
\begin{ex}
% ID: L12C1B4-94-TLN
% Mức: NB
Cùng một quá trình, tăng công suất từ 800 $W$ lên 1200 $W$, $H$ không đổi. Thời gian mới bằng bao nhiêu phần trăm thời gian cũ? Chỉ ghi số.
\shortans{66.67}
\loigiai{
$t2/t1=P1/P2=800/1200=2/3=66,67%.$
}
\end{ex}

% ===== Câu 8 =====
% ID: L12C1B4-95-TN
% Mức: NB
\begin{ex}
% ID: L12C1B4-95-TN
% Mức: NB
Nguồn công suất 2000 $W$, hiệu suất 75%, hóa hơi $0,4\,\text{kg}$ chất có L=1,5.10^ $6\,\text{J}$ /kg ở nhiệt độ sôi. Thời gian là
\choice
{$40\,\text{s}$}
{\True $400\,\text{s}$}
{$4000\,\text{s}$}
{400 phút}
\loigiai{
H.P.t=Lm nên t=Lm/(HP)= $400\,\text{s}$.
}
\end{ex}

% ===== Câu 9 =====
% ID: L12C1B4-96-DS
% Mức: VD
\begin{ex}
% ID: L12C1B4-96-DS
% Mức: VD
Nguồn 1000 $W$, hiệu suất 80%, hoạt động $300\,\text{s}$.
\choiceTF
{\True Điện năng tiêu thụ là 300 kJ.}
{\True Năng lượng hữu ích là 240 kJ.}
{\True Nếu L=2,0.10^ $6\,\text{J}$ /kg thì hóa hơi $0,12\,\text{kg}$.}
{Hiệu suất càng nhỏ thì cùng $P$,t hóa hơi càng nhiều.}
\loigiai{
\begin{itemchoice}
\item (Đúng) Điện năng tiêu thụ là 300 kJ.
\item (Đúng) Năng lượng hữu ích là 240 kJ.
\item (Đúng) Nếu L=2,0.10^ $6\,\text{J}$ /kg thì hóa hơi $0,12\,\text{kg}$.
\item (Sai) Hiệu suất càng nhỏ thì cùng $P$,t hóa hơi càng nhiều.
\end{itemchoice}
}
\end{ex}

% ===== Câu 10 =====
% ID: L12C1B4-97-TL
% Mức: TH
\begin{bt}
% ID: L12C1B4-97-TL
% Mức: TH
Một ấm vừa đun nóng nước đến sôi vừa làm bay hơi. Viết công thức thời gian có kể nhiệt dung ấm và hiệu suất.
\loigiai{
Nhiệt hữu ích Q=(mc+C_ấm) $(t_s-t_0)$ +Lm_h. Vì Q=HPt nên t=[(mc+C_ấm) $(t_s-t_0)$ +Lm_h]/(HP).
}
\end{bt}

% ===== Câu 11 =====
% ID: L12C1B4-98-TLN
% Mức: TH
\begin{ex}
% ID: L12C1B4-98-TLN
% Mức: TH
Hóa hơi $0,20\,\text{kg}$ chất có L=1,8.10^ $6\,\text{J}$ /kg trong $400\,\text{s}$ bằng nguồn $1200\,\text{W}$. Tính hiệu suất theo phần trăm, chỉ ghi số.
\shortans{75}
\loigiai{
$H=Lm/(Pt)=1,8.10^6.0,20/(1200.400)=0,75=75%.$
}
\end{ex}

% ===== Câu 12 =====
% ID: L12C1B4-99-TN
% Mức: NB
\begin{ex}
% ID: L12C1B4-99-TN
% Mức: NB
Muốn giảm một nửa thời gian hóa hơi cùng khối lượng khi $H$ không đổi, cần
\choice
{giảm $P$ một nửa}
{\True tăng $P$ gấp đôi}
{giữ $P$ không đổi}
{tăng $L$ gấp đôi}
\loigiai{
t tỉ lệ nghịch với P.
}
\end{ex}

% ===== Câu 13 =====
% ID: L12C1B4-100-DS
% Mức: VD
\begin{ex}
% ID: L12C1B4-100-DS
% Mức: VD
Về xác định hiệu suất khi hóa hơi.
\choiceTF
{\True $H=Lm/(Pt).$}
{\True $H$ không vượt quá 1 trong mô hình thông thường.}
{\True Nếu bỏ qua hao phí thì H=1.}
{H=Pt/(Lm) luôn đúng.}
\loigiai{
\begin{itemchoice}
\item (Đúng) $H=Lm/(Pt).$.
\item (Đúng) $H$ không vượt quá 1 trong mô hình thông thường.
\item (Đúng) Nếu bỏ qua hao phí thì H=1.
\item (Sai) H=Pt/(Lm) luôn đúng.
\end{itemchoice}
}
\end{ex}

% ===== Câu 14 =====
% ID: L12C1B4-101-TL
% Mức: VD
\begin{bt}
% ID: L12C1B4-101-TL
% Mức: VD
Một thiết bị 1000 $W$ chạy 8 phút, làm hóa hơi $0,15\,\text{kg}$ chất có L=2,0.10^ $6\,\text{J}$ /kg. Xác định hiệu suất.
\loigiai{
$H=Lm/(Pt)=2,0.10^6.0,15/(1000.480)=0,625=62,5%.$
}
\end{bt}

% ===== Câu 15 =====
% ID: L12C1B4-102-TLN
% Mức: VD
\begin{ex}
% ID: L12C1B4-102-TLN
% Mức: VD
Thiết bị hiệu suất 80% hóa hơi $0,30\,\text{kg}$ chất có L=2,0.10^ $6\,\text{J}$ /kg trong $500\,\text{s}$. Tính công suất theo $W$, chỉ ghi số.
\shortans{1500}
\loigiai{
$P=Lm/(Ht)=2,0.10^6.0,30/(0,8.500)=1500 W.$
}
\end{ex}

% ===== Câu 16 =====
% ID: L12C1B4-103-TN
% Mức: TH
\begin{ex}
% ID: L12C1B4-103-TN
% Mức: TH
Nguồn công suất 1200 $W$, hiệu suất 75%, hóa hơi $0,15\,\text{kg}$ chất có L=1,8.10^ $6\,\text{J}$ /kg ở nhiệt độ sôi. Thời gian là
\choice
{$30\,\text{s}$}
{\True $300\,\text{s}$}
{$3000\,\text{s}$}
{300 phút}
\loigiai{
H.P.t=Lm nên t=Lm/(HP)= $300\,\text{s}$.
}
\end{ex}

% ===== Câu 17 =====
% ID: L12C1B4-104-DS
% Mức: VDC
\begin{ex}
% ID: L12C1B4-104-DS
% Mức: VDC
Xét công thức HPt=Lm.
\choiceTF
{\True $H$ phải viết dạng thập phân.}
{\True t tính bằng giây nếu $P$ tính W.}
{\True $P$ tính $W$ tức là J/s.}
{Có thể dùng phút trực tiếp mà không đổi.}
\loigiai{
\begin{itemchoice}
\item (Đúng) $H$ phải viết dạng thập phân.
\item (Đúng) t tính bằng giây nếu $P$ tính W.
\item (Đúng) $P$ tính $W$ tức là J/s.
\item (Sai) Có thể dùng phút trực tiếp mà không đổi.
\end{itemchoice}
}
\end{ex}

% ===== Câu 18 =====
% ID: L12C1B4-105-TL
% Mức: VD
\begin{bt}
% ID: L12C1B4-105-TL
% Mức: VD
Phân tích vì sao không thể lấy Pt=Lm khi hiệu suất chỉ 70%.
\loigiai{
Pt là toàn bộ năng lượng nguồn; chỉ phần HPt=0,70Pt truyền hữu ích cho chất. Phần còn lại thất thoát. Vì vậy phương trình đúng là 0,70Pt=Lm, nếu chất đã ở nhiệt độ sôi.
}
\end{bt}

% ===== Câu 19 =====
% ID: L12C1B4-106-TLN
% Mức: VDC
\begin{ex}
% ID: L12C1B4-106-TLN
% Mức: VDC
Nguồn 1200 $W$, hiệu suất 75%, hóa hơi $0,18\,\text{kg}$ chất có L=1,5.10^ $6\,\text{J}$ /kg. Tính thời gian theo giây, chỉ ghi số.
\shortans{300}
\loigiai{
$t=Lm/(HP)=1,5.10^6.0,18/(0,75.1200)=300 s.$
}
\end{ex}

% ===== Câu 20 =====
% ID: L12C1B4-107-TN
% Mức: TH
\begin{ex}
% ID: L12C1B4-107-TN
% Mức: TH
Hệ thức đúng khi nguồn công suất $P$, hiệu suất $H$ chỉ dùng để hóa hơi khối lượng m là
\choice
{$Pt=HLm$}
{\True $HPt=Lm$}
{$Pt=Lm/H^2$}
{$Ht=PLm$}
\loigiai{
Năng lượng hữu ích HPt bằng nhiệt lượng hóa hơi Lm.
}
\end{ex}

% ===== Câu 22 =====
% ID: L12C1B4-108-TL
% Mức: VDC
\begin{bt}
% ID: L12C1B4-108-TL
% Mức: VDC
Nguồn có hiệu suất 75% phải hóa hơi $0,50\,\text{kg}$ chất trong 10 phút, L=1,8.10^ $6\,\text{J}$ /kg. Tìm công suất tối thiểu.
\loigiai{
$P=Lm/(Ht)=1,8.10^6.0,50/(0,75.600)=2000 W.$
}
\end{bt}

% ===== Câu 23 =====
% ID: L12C1B4-109-TLN
% Mức: VDC
\begin{ex}
% ID: L12C1B4-109-TLN
% Mức: VDC
Nguồn 2000 $W$, hiệu suất 60%, hoạt động 5 phút. Năng lượng hữu ích theo kJ là bao nhiêu? Chỉ ghi số.
\shortans{360}
\loigiai{
$Q_h=0,6.2000.300=360000 J=360 kJ.$
}
\end{ex}

% ===== Câu 24 =====
% ID: L12C1B4-110-TN
% Mức: TH
\begin{ex}
% ID: L12C1B4-110-TN
% Mức: TH
Nguồn 1000 $W$ hoạt động 5 phút, hiệu suất 60%. Năng lượng hữu ích là
\choice
{\True 180 kJ}
{300 kJ}
{500 kJ}
{600 kJ}
\loigiai{
$Q_h=HPt=0,60.1000.300=180000 J=180 kJ.$
}
\end{ex}

% ===== Câu 28 =====
% ID: L12C1B4-111-TN
% Mức: VD
\begin{ex}
% ID: L12C1B4-111-TN
% Mức: VD
Nguồn công suất 800 $W$, hiệu suất 90%, hóa hơi $0,09\,\text{kg}$ chất có L=2,4.10^ $6\,\text{J}$ /kg ở nhiệt độ sôi. Thời gian là
\choice
{$30\,\text{s}$}
{\True $300\,\text{s}$}
{$3000\,\text{s}$}
{300 phút}
\loigiai{
H.P.t=Lm nên t=Lm/(HP)= $300\,\text{s}$.
}
\end{ex}

% ===== Câu 32 =====
% ID: L12C1B4-112-TN
% Mức: VD
\begin{ex}
% ID: L12C1B4-112-TN
% Mức: VD
Nếu $P$ và $H$ không đổi, thời gian hóa hơi t tỉ lệ
\choice
{nghịch với m}
{\True thuận với m}
{với m^2}
{không phụ thuộc m}
\loigiai{
t=Lm/(HP), nên t tỉ lệ thuận m.
}
\end{ex}

% ===== Câu 36 =====
% ID: L12C1B4-113-TN
% Mức: VDC
\begin{ex}
% ID: L12C1B4-113-TN
% Mức: VDC
Nguồn công suất 1500 $W$, hiệu suất 80%, hóa hơi $0,3\,\text{kg}$ chất có L=1,2.10^ $6\,\text{J}$ /kg ở nhiệt độ sôi. Thời gian là
\choice
{$30\,\text{s}$}
{\True $300\,\text{s}$}
{$3000\,\text{s}$}
{300 phút}
\loigiai{
H.P.t=Lm nên t=Lm/(HP)= $300\,\text{s}$.
}
\end{ex}

% ===== Câu 40 =====
% ID: L12C1B4-114-TN
% Mức: VDC
\begin{ex}
% ID: L12C1B4-114-TN
% Mức: VDC
Hiệu suất 80% được thay vào công thức dưới dạng
\choice
{80}
{8}
{\True 0,8}
{0,08}
\loigiai{
$80%=0,80.$
}
\end{ex}
